\documentclass[12pt,a4paper]{article}
\usepackage[english]{babel}
\usepackage[T1]{fontenc}
\usepackage[margin=18mm]{geometry}
\usepackage{amsmath,amssymb}
\usepackage{enumitem}
\usepackage[most]{tcolorbox}
\usepackage[colorlinks=true,linkcolor=blue]{hyperref}
\setlength{\parindent}{0pt}
\setlist{topsep=0.5\baselineskip,itemsep=0.5\baselineskip,parsep=0pt,leftmargin=2.45em}
\newcounter{definition}
\setcounter{definition}{4}
\renewcommand{\thedefinition}{3.\arabic{definition}}
\definecolor{NotesDefinition}{rgb}{0.70,0.94,0.70}
\definecolor{NotesPaper}{gray}{0.95}
\tcbset{notesbox/.style={enhanced,breakable,
  colback=NotesPaper,coltitle=black,fonttitle=\bfseries,
  boxrule=0.5mm,arc=1mm,boxsep=1mm,left=1mm,right=1mm,top=1mm,bottom=1mm,
  before skip=7mm,after skip=4mm,
  attach boxed title to top left={xshift=12.5mm,yshift=-0.5mm},
  boxed title style={enhanced,boxrule=0.4mm,arc=0.75mm,
    boxsep=1mm,left=2mm,right=2mm,top=0mm,bottom=0mm,
    frame code={%
      \path[fill=tcbcolframe,draw=tcbcolframe]
        (frame.north west) -- (frame.north east)
        .. controls ([xshift=4mm]frame.north east) and
                    ([xshift=7mm]frame.south east) ..
                    ([xshift=11mm]frame.south east)
        -- ([xshift=-11mm]frame.south west)
        .. controls ([xshift=-7mm]frame.south west) and
                    ([xshift=-4mm]frame.north west) .. (frame.north west) -- cycle;
    }},
  before upper={\setlength{\parskip}{0.5\baselineskip}}},
  definitionbox/.style={notesbox,colframe=NotesDefinition,
    boxed title style={colback=NotesDefinition,colframe=NotesDefinition}}}
\newtcolorbox[use counter=definition]{frameddefn}[1]{definitionbox,
  title={Definition \thedefinition:\enspace #1}}
\renewcommand{\thedefinition}{3.\arabic{definition}}
\begin{document}
\begin{center}
  {\Large\bfseries Fundamentals of Statistical Learning}\par
  \smallskip
  {\large Definitions cheatsheet}
\end{center}

\setcounter{definition}{1}
\begin{frameddefn}{Random variable}\label{def:random-variable}
A \emph{random variable} is a measurable map
\[
  X:\Omega\longrightarrow\mathbb R.
\]
It assigns one real number $X(\omega)$ to each outcome $\omega$.
\end{frameddefn}

\setcounter{definition}{3}
\begin{frameddefn}{Preimage}\label{def:preimage}
For a random variable $X:\Omega\to\mathbb R$ and a set
$E\subseteq\mathbb R$, the \emph{preimage} (or \emph{inverse image}) of
$E$ is
\[
  X^{-1}(E):=\{\omega\in\Omega:X(\omega)\in E\}\subseteq\Omega.
\]
To find this event, we collect the outcomes whose outputs fall in $E$.
Whenever $X^{-1}(E)$ is an event, we define
\[
  \begin{aligned}
    P(X\in E)&:=P\bigl(X^{-1}(E)\bigr)\\
              &=P\bigl(\{\omega\in\Omega:X(\omega)\in E\}\bigr).
  \end{aligned}
\]
In particular, for a single value $x$,
\[
  P(X=x):=P\bigl(X^{-1}(\{x\})\bigr).
\]
Here $X$ is the random variable, while $x$ is a particular real number.
\end{frameddefn}

\setcounter{definition}{4}
\begin{frameddefn}{Distribution of a random variable}
\label{def:distribution}
The \emph{distribution} of $X$ is the rule that tells us with what
probability $X$ takes values in a given set $E$.
\[
  P_X(E)=P(X\in E).
\]
\end{frameddefn}

\newpage
\setcounter{definition}{5}
\begin{frameddefn}{Discrete random variable and PMF}
\label{def:discrete-pmf}
A random variable $X$ is \emph{discrete} if there is a finite or
countably infinite set\begingroup\renewcommand{\thefootnote}{\fnsymbol{footnote}}\footnotemark[1]\endgroup
$\mathcal X=\{x_1,x_2,\ldots\}$ such that
$P(X\in\mathcal X)=1$.

Its \emph{probability mass function} (PMF) is
\[
  p_X(x):=P(X=x),\qquad x\in\mathbb R.
\]
\begin{enumerate}
  \item $p_X(x)\geq0$ for every $x\in\mathbb R$.
  \item The \emph{support}, in the discrete sense used here, is
    $\operatorname{supp}(X):=\{x\in\mathbb R:p_X(x)>0\}$.
    Thus the support lists the values that have positive probability;
    values outside it have probability zero.
  \item $\displaystyle\sum_{x\in\mathcal X}p_X(x)=1$.
\end{enumerate}
\end{frameddefn}

\begingroup
\renewcommand{\thefootnote}{\fnsymbol{footnote}}
\footnotetext[1]{A finite set has finitely many elements. A set $A$ is countably infinite if there exists a bijection $f:\mathbb N\to A$.}
\endgroup

\setcounter{definition}{6}
\begin{frameddefn}{Continuous random variable and PDF}
\label{def:continuous-pdf}
A random variable $X$ is \emph{continuous} (in the sense used here)
if there exists a function $f_X$, called its \emph{probability density
function} (PDF), such that, for every $a\leq b$,
\[
  P(a\leq X\leq b)=\int_a^b f_X(x)\,dx.
\]
\begin{enumerate}
  \item $f_X(x)\geq0$ for every $x\in\mathbb R$.
  \item Following the convention in the slides, the \emph{support} is
    $\operatorname{supp}(X):=\{x\in\mathbb R:f_X(x)>0\}$.
  \item $\displaystyle\int_{\mathbb R}f_X(x)\,dx=1$.
\end{enumerate}
\end{frameddefn}

\setcounter{definition}{7}
\begin{frameddefn}{Cumulative distribution function}
\label{def:cdf}
Given a random variable $X$, its \emph{cumulative distribution
function} (CDF) is the function $F_X:\mathbb R\to[0,1]$ defined by
\[
  F_X(x):=P(X\leq x)=P\bigl(X\in(-\infty,x]\bigr),
  \qquad x\in\mathbb R.
\]
It gives the probability accumulated up to and including the
threshold $x$. It is defined for both
\hyperref[def:discrete-pmf]{discrete} and
\hyperref[def:continuous-pdf]{continuous} random variables.
\end{frameddefn}

\setcounter{definition}{9}
\begin{frameddefn}{Quantile function: general case}
\label{def:quantile-general}
Let $X$ be any random variable, discrete or continuous, with CDF
$F_X$. For a probability level $0<p<1$, its \emph{quantile function} is
\[
  Q(p):=\inf\{x\in\mathbb R:F_X(x)\geq p\}.
\]
It gives the smallest threshold $x$ such that
$P(X\leq x)$ is at least $p$. We also write $q_p$ instead of $Q(p)$.
\end{frameddefn}

\clearpage
\section*{Relations}

\subsection*{From the PMF to the CDF}
\begin{tcolorbox}[colback=blue!3,colframe=blue!30!black,
  boxrule=0.4pt,arc=1mm,title={What we are trying to answer},
  colbacktitle=blue!10,coltitle=black,fonttitle=\bfseries]
If we know the probability of each individual
value of $X$, how do we calculate the probability that $X$ is at most $x$?
\end{tcolorbox}

For a discrete random variable, the event $X\leq x$ includes every
possible value $k$ that is at most $x$. These events $\{X=k\}$ are
pairwise disjoint, so we add their probabilities:
\[
  \boxed{\underbrace{F_X(x)}_{\text{CDF: to find}}=P(X\leq x)
    =\sum_{\substack{k\in\operatorname{supp}(X)\\k\leq x}}
      \underbrace{p_X(k)}_{\text{PMF: known}},
    \qquad x\in\mathbb R.}
\]
\subsection*{From the PDF to the CDF}
\label{rel:pdf-to-cdf}
\begin{tcolorbox}[colback=blue!3,colframe=blue!30!black,
  boxrule=0.4pt,arc=1mm,title={What we are trying to answer},
  colbacktitle=blue!10,coltitle=black,fonttitle=\bfseries]
We know how probability is spread over the possible values of $X$
(its PDF). How do we calculate the probability of getting a value
less than or equal to $x$?
\end{tcolorbox}
For a continuous random variable with PDF $f_X$, we calculate
$P(X\leq x)$ by integrating the PDF over $(-\infty,x]$:
\[
  \boxed{\underbrace{F_X(x)}_{\text{CDF: to find}}=P(X\leq x)
    =\int_{-\infty}^{x}\underbrace{f_X(t)}_{\text{PDF: known}}\,dt,
    \qquad x\in\mathbb R.}
\]
\clearpage
\subsection*{From the CDF to the PMF and PDF}
\begin{tcolorbox}[colback=blue!3,colframe=blue!30!black,
  boxrule=0.4pt,arc=1mm,title={What we are trying to answer},
  colbacktitle=blue!10,coltitle=black,fonttitle=\bfseries]
We know the CDF $F_X$. How do we recover the PMF $p_X$ if $X$ is
discrete, or the PDF $f_X$ if $X$ is continuous?
\end{tcolorbox}

\textbf{What we know.} The CDF is defined by
\[
  F_X(x):=P(X\leq x).
\]
For every threshold $x$, it gives the probability of obtaining a
value less than or equal to $x$. We want to recover the PMF or the
PDF, depending on the type of $X$.

\textbf{Discrete case: recover the PMF.}
The PMF is defined by
\[
  p_X(x):=P(X=x).
\]
The CDF includes all values up to $x$. To isolate the probability
of the single value $x$, we subtract the probability of values
strictly less than $x$:
\[
  P(X=x)=P(X\leq x)-P(X<x).
\]
Since the left limit satisfies
\[
  F_X(x^-):=\lim_{t\uparrow x}F_X(t)=P(X<x),
\]
we obtain
\[
  \boxed{\underbrace{p_X(x)}_{\text{PMF: to find}}
    =\underbrace{F_X(x)-F_X(x^-)}_{\text{CDF: known}}.}
\]
Thus the PMF at $x$ is the size of the jump of the CDF at $x$.

\textbf{Continuous case: recover the PDF.}
A PDF is a non-negative function $f_X$ such that
\[
  P(a\leq X\leq b)=\int_a^b f_X(t)\,dt,
  \qquad \int_{\mathbb R}f_X(t)\,dt=1.
\]
In particular,
\[
  F_X(x)=P(X\leq x)=\int_{-\infty}^{x}f_X(t)\,dt.
\]
The CDF is therefore the integral of the PDF. To recover a PDF,
we differentiate the CDF:
\[
  \boxed{\underbrace{f_X(x)}_{\text{PDF: to find}}
    =\frac{d}{dx}\underbrace{F_X(x)}_{\text{CDF: known}}.}
\]
This recovers a PDF at points where the derivative exists, by the
fundamental theorem of calculus.


\end{document}
